\chapter*{Conclusion générale}
\addcontentsline{toc}{chapter}{Conclusion générale}

Ce stage d'un mois chez Forvis Mazars Group m'a permis de réaliser un prototype complet autour de la transformation de renseignements CTI en propositions de détection. Ce travail m'a amené à combiner des compétences de développement backend, frontend, base de données, orchestration asynchrone, cybersécurité et documentation technique.

La plateforme obtenue intègre MISP, normalise des événements CTI, exécute un workflow LangGraph, utilise une abstraction IA, génère des candidats Sigma, applique des watchers, valide avec pySigma, détecte les doublons et soumet les propositions à une revue humaine. Les pages frontend permettent de visualiser la file des menaces, les sessions IA, les runs de graphe, les propositions, la matrice ATT\&CK, la télémétrie, la santé système et le catalogue.

Le principal apprentissage du stage est que l'intelligence artificielle ne peut pas être ajoutée seule à un processus de cybersécurité. Elle doit être entourée de validations déterministes, d'une traçabilité forte, d'une mémoire structurée, de règles de sécurité et d'un contrôle humain. Le projet illustre cette approche en séparant la génération IA de la décision finale.

Le travail reste néanmoins un prototype académique. Pour une utilisation en production, il faudrait renforcer l'authentification, ajouter une observabilité complète, automatiser les tests frontend, valider le mode IA réel dans des conditions stables, élargir la couverture ATT\&CK et connecter la publication à un SIEM réel. Ces perspectives constituent une continuité naturelle du travail réalisé.
